\section{Monte Carlo with Exploring Starts}
    \subsection{Main Idea}
        The MC-Basic algorithms is not sample-efficient, since it only updates only q value of
        one state-action pair each episode. To deal with this low efficiency, we use all the
        \emph{visit} to the state-action pair in each episode.
        
        Suppose that the agent has gone through following trajectory:
        \begin{equation}
            s_1 \xrightarrow{a_1} s_2 \xrightarrow{a_2} s_3 \xrightarrow{a_3} s_4 \xrightarrow{a_4} \dots
        \end{equation}
        To increase the sample efficiency, we use the visits to the state-action pairs in the trajectory.
        That is, using the above episode, we not only use $(s_1,a_1)$ for estimating $q(s_1,a_1)$ but use $(s_2,a_2)$ for estimating $q(s_2,a_2)$, and so on.
        There are two different ways to do this:
        \begin{itemize}
            \item The first way is to only use the state-action pairs when it is first encountered. This is called \emph{first-visit method};
            \item The second way is to use all the visit to the state-action pairs. This is called \emph{every-visit method}.
        \end{itemize}

        Also, we update the estimate of q values after each episode ends. This falls into the scope of \emph{generalized policy iteration},
        since we are using the value approximation even it is not so accurate.
    \subsection{Pseudocode}
        \Algorithm{mc-explore}
        {Monte Carlo with Exploring Starts (every-visit)}
        {
            \inputitem{$\pi_0$}{initial policy}
        }
        {
            \inputitem{$\pi^*$}{optimal policy}
        }
        {
            \tcp{initialize}
            \ForEach{valid state-action pair $(s,a)$}
            {
                $\pi \leftarrow \pi_0$\;
                $q(s,a) \leftarrow$ initial value\;
                \tcp{holds the sum of returns of (sub)episode starting from $(s,a)$}
                $\text{Ret}(s,a) \leftarrow 0$\;
                \tcp{counts the number of visits of $(s,a)$}
                $\text{N}(s,a) \leftarrow 0$\;
            }
            \Repeat{convergence}
            {
                select a starting valid state-action pair $(s,a)$, ensuring that all the pairs can be possibly selected\;
                generate an episode starting from $(s,a)$, following current policy $\pi$: $s_0, a_0, r_1, s_1, a_1, \dots, s_{T-1}, a_{T-1}, r_T$\;
                $G \leftarrow 0$\;
                \For{$t=T-1,T-2,\dots,0$}
                {
                    $G \leftarrow r_{t+1} + \gamma G$\;
                    $\text{Ret}(s,a) \leftarrow \text{Ret}(s,a) + G$\;
                    $\text{N}(s,a) \leftarrow \text{N}(s,a) + 1$\;
                    \tcp{policy evaluation}
                    $q(s_t,a_t) \leftarrow \frac{\text{Ret}(s,a)}{\text{N}(s,a)}$\;
                    \tcp{policy improvement}
                    $\pi_{k+1}(a|s) = \begin{cases}
                        1 &\text{if $a=\arg\max_{a'}{q_{\pi_k}(s,a')}$} \\
                        0 &\text{otherwise}
                    \end{cases}$\;
                }
            }
        }
        
    \subsection{Implementation Details}
        The exploring-start condition can be achieved by random start state.

    \subsection{Pros and Cons}
    Pros:
    \begin{itemize}
        \item Better sample efficiency then MC-Basic;
        \item All state-action pairs can be visited due to the exploring-start, and this makes the estimation more accurate.
    \end{itemize}
    Cons:
    \begin{itemize}
        \item Only applicable to episodic tasks, which is a common drawback for all MC methods. This property leads to following problem, too;
        \item High variance since the return holds all the noises and randomnesses, and the return is used for update;
        \item If the exploring start condition cannot be satisfied, the MC with Exploring Starts cannot be used. This is common in real-world problems.
        To deal with this, another algorithms without exploring starts will be introduced next section;
    \end{itemize}